%% Proof annex fragment for Part 2.2 normal forms.  R-007 and R-008 are
%% retained in the individual constructions rather than hidden in citations.

\begin{annexitem}{annex:p2.fol.nf.1}{thm:p02-fol-nnf}{First-order negation normal form}
\begin{proof}
Define the two maps simultaneously.  They send an atom \(A\) to
\(A,\neg A\), and \(\bottom\) to \(\bottom,\neg\bottom\), respectively;
they exchange under an outer negation.  The positive map commutes with
\(\land,\lor\), replaces \(\varphi\to\psi\) by
\(\operatorname{nnf}^{\neg}(\varphi)\lor\operatorname{nnf}(\psi)\), and
commutes with either quantifier.  The companion swaps \(\land/\lor\), sends
the implication to
\(\operatorname{nnf}(\varphi)\land\operatorname{nnf}^{\neg}(\psi)\), and
exchanges \(\forall/\exists\).  Calls have smaller complexity.

Simultaneous strong induction proves the outputs are generated by the NNF
constructors.  The satisfaction clauses give equivalence, including the two
quantifier dualities under every updated assignment.  The same recursion
neither creates nor deletes variable occurrences, so the recursive
\(\FV\) and \(\operatorname{Var}\) clauses give both asserted equalities.
\end{proof}
\end{annexitem}

\begin{annexitem}{annex:p2.fol.nf.2}{thm:p02-fol-rectification}{Rectified variant}
\begin{proof}
Define \(\operatorname{rect}\) recursively under
\(\operatorname{Var}(\varphi)\subseteq V\).  Atoms and \(\bottom\) are
fixed; commute with negation.  For
\(\psi\circ\chi\), first compute \(\operatorname{rect}(\psi,V)\), then compute the
right side with every variable of that output added to \(V\).  For
\(Qx\psi\), choose the least \(z\notin V\) and put
\[
 \operatorname{rect}(Qx\psi,V)
 =Qz\,\operatorname{rect}(\psi[z/x],V\cup\{z\}).
\]
Because \(z\notin\operatorname{Var}(\psi)\), substitution reaches no
binder named \(z\), is textual renaming in atomic formulas, preserves the
invariant, and has complexity
\(c(\psi[z/x])=c(\psi)<c(Qx\psi)\).  Thus this is a legitimate
strong-complexity recursion, not a subformula induction on a renamed body.

Alpha-equivalence preserves \(\FV\): equality of free-variable sets is a
constructor congruence containing every generating safe-renaming pair.
Strong induction now gives
\(\operatorname{rect}(\varphi,V)=_{\alpha}\varphi\); the quantifier step combines the
induction hypothesis, quantifier congruence, and the safe-renaming generator.
Each chosen binder avoids \(V\); the binary clause makes right binders avoid
left-output variables, and recursion makes inner binders avoid outer ones.
Since alpha-equivalence preserves \(\FV\subseteq V\), no binder is free.
Finally every clause preserves the constructor tree; quantifier renaming
changes variables only inside atoms.  Hence rectification, free variables,
and NNF are all preserved.
\end{proof}
\end{annexitem}

\begin{annexitem}{annex:p2.fol.nf.3}{thm:p02-fol-prenex-extraction}{Prenex extraction}
\begin{proof}
First, if \(y\notin\FV(\psi)\), satisfaction and coincidence give
\[
 (Qy\varphi)\circ\psi\Leftrightarrow Qy(\varphi\circ\psi),
 \qquad
 \psi\circ(Qy\varphi)\Leftrightarrow Qy(\psi\circ\varphi)
\]
for \(Q\in\{\forall,\exists\}\) and
\(\circ\in\{\land,\lor\}\).  For example, if
\(\psi\) is true, both \((\forall y\,\varphi)\lor\psi\) and
\(\forall y(\varphi\lor\psi)\) are true.  If \(\psi\) is false,
coincidence makes it false under every \(y\)-update, so both sides reduce to
\(\forall y\,\varphi\).  The other seven cases are the same satisfaction
calculation.  Formula equivalence is also a
congruence for both connectives and quantifiers.

Induct on the NNF generation.  Literals and units are fixed; at \(Qx\varphi\)
prefix \(Qx\) to the inductive output.  For
\(\varphi\circ\psi\), write the two outputs as
\(Q_1y_1\cdots Q_ky_k\varphi_0\) and
\(R_1z_1\cdots R_mz_m\psi_0\), and concatenate their prefixes before
\(\varphi_0\circ\psi_0\).  Rectification makes each \(y_i\) absent from
\(\FV(\psi)\) and each \(z_j\) absent from
\(\FV(\varphi_0)\) and distinct from every \(y_i\).  Repeatedly applying the
displayed laws therefore moves exactly those existing quantifiers outward.
Congruence and induction give equivalence.  No variable or quantifier is
introduced, deleted, or renamed, so the free/all-variable sets, binder
distinctness, exact quantifier occurrences, and quantifier-free NNF matrix
have the stated values.
\end{proof}
\end{annexitem}

\begin{annexitem}{annex:p2.fol.nf.4}{thm:p02-fol-prenex}{Prenex normal form}
\begin{proof}
Apply Theorem~\ref{thm:p02-fol-nnf}, rectify with forbidden set
\(\operatorname{Var}(\operatorname{nnf}(\varphi))\), then use
Theorem~\ref{thm:p02-fol-prenex-extraction}.  Each operation stays in
\(\mathcal L\), preserves equivalence, and supplies the next hypotheses.
Transitivity gives the required rectified prenex equivalent of \(\varphi\).
\end{proof}
\end{annexitem}

\begin{annexitem}{annex:p2.fol.nf.5}{thm:p02-skolem-equisatisfiable}{Expansion--reduct theorem for Skolemization}
\begin{proof}
Write \(\sigma=Q_1x_1\cdots Q_nx_n\chi\), and let \(\sigma_i\) be the
suffix beginning after position \(i\).  Let \(q(i)\) count preceding
universal positions.  At existential position \(j\), write
\(t_j=f_j(\bar u(j))\) for the Skolem term in the preceding universal
variables.  Then
\(\operatorname{Sk}(\sigma)=\forall\bar u\,\chi^*\), where
\(\chi^*\) is the simultaneous replacement of every existential
\(x_j\) by \(t_j\).  The matrix is quantifier-free and no \(t_j\) contains
an existential variable, so these replacements may be composed in any
order.

Suppose \(\mathcal M\entails\sigma\).  At each existential position choose
recursively
\[
 F_j:(\dom\mathcal M)^{q(j-1)}\longrightarrow\dom\mathcal M
\]
while maintaining
\[
 \mathcal M,s_{\bar a}\entails\sigma_i
 \quad\text{for every }\bar a\in(\dom\mathcal M)^{q(i)}. \tag{I}
\]
Here \(s_{\bar a}\) assigns preceding universal variables from \(\bar a\)
and each earlier existential variable the value given by its \(F_j\).
The universal step appends an arbitrary domain element.  At an existential
step, (I) gives for every \(\bar a\) a nonempty witness set; AC chooses one
member uniformly to define \(F_{i+1}\).  This is the only use of choice.

Interpret each new \(f_j\) by \(F_j\).  At \(i=n\), invariant (I), reduct
invariance, term evaluation, and repeated use of the substitution lemma give
the expanded structure satisfaction of \(\chi^*\) under every tuple of
universal values.  Hence it satisfies \(\operatorname{Sk}(\sigma)\).

Conversely let \(\mathcal N\entails\operatorname{Sk}(\sigma)\).  Now every
existential value is already the value of its displayed Skolem term.  Downward
induction on \(i\) proves
\[
 \mathcal N,s_{\bar a}\entails\sigma_i
 \quad\text{for every }\bar a\in(\dom\mathcal N)^{q(i)}. \tag{II}
\]
The base reverses the same iterated substitution computation from
\(\chi^*\) to \(\chi\).  A universal step uses (II) for every appended
value; an existential step uses the already interpreted Skolem-term value as
its witness.  Thus (II) at zero gives \(\mathcal N\entails\sigma\), and
reduct invariance yields
\(\mathcal N\!\upharpoonright\mathcal L\entails\sigma\).  No selection is
made in this direction.
\end{proof}
\end{annexitem}
